\section{Introduction}
\label{sec:introduction}

\subsection{The Local-Optimum Problem in METIS Bisection}

The core computational primitive of the B-series redistricting platform is
the bisection node: given a subgraph of $m$ census tracts that must be
divided into two contiguous, population-balanced pieces, find the cut that
minimises the weighted edge cut~\citep{b8paper, b11paper}.
The standard solver is \metis{}, a multilevel graph partitioner that
coarsens the graph, finds an approximate minimum bisection on the coarsened
graph, and uncoarsens with Kernighan-Lin (KL) local refinement~\citep{karypis1998}.
\metis{} is fast---$O(m \log m)$ per call---and produces solutions within a
small constant factor of optimal on typical redistricting graphs.

However, \metis{} is a heuristic.
Graph minimum bisection is $\mathcal{NP}$-hard~\citep{garey1976}, and the KL
refinement phase terminates at a local optimum that depends on the random
seed.
The B.7 analysis showed that for most states, seed variation has little
practical effect on the final map, but the variation is nonzero and
occasionally significant~\citep{b7paper}.
The B-series response to this fragility has followed two tracks.

\paragraph{Track 1: Better seed search.}
\cs{} (B.16) sweeps forward through seeds, running \metis{} at each, and
halts when $T=600$ consecutive seeds produce no improvement in normalised
edge cut~\citep{b16paper}.
\be{} (H.1) runs $T$ independent \metis{} seeds in parallel and selects
the minimum~\citep{h1paper}.
Both are \emph{SeedCompositor} strategies---they change how many times the
structure-layer bisector is called, not how it works internally.

\paragraph{Track 2: Better structure-layer bisection.}
\geosec{} (B.8) replaces the flat tree with a ratio-optimal tree that
normalises edge cut by $\sqrt{\min(i,k-i)}$~\citep{b8paper}.
\areasec{} (B.9) adds an area-swing constraint to bias toward
geographically compact pieces~\citep{b9paper}.
Both change the \emph{objective} at each node while still delegating
the solver to \metis{}.

This paper introduces a third track: replace \metis{} itself at each node
with a \textbf{simulated annealing} loop that refines the \metis{} solution
by accepting worse proposals probabilistically, guided by a geometric
cooling schedule.

\subsection{Simulated Annealing as a Bisection Solver}

Simulated Annealing (SA) is a classical meta-heuristic for combinatorial
optimisation~\citep{kirkpatrick1983}.
The algorithm starts at an initial solution and iteratively proposes
small perturbations.
Improvements are always accepted.
Degradations are accepted with probability $\exp(-\Delta f / T)$, where
$\Delta f > 0$ is the cost increase and $T > 0$ is a temperature parameter
that decreases over time according to a \emph{cooling schedule}.
At high temperature the algorithm accepts most proposals (exploration);
as $T \to 0$ it becomes a greedy local-search descent (exploitation).
In the limit of infinitely slow cooling, SA converges to the global
optimum with probability one~\citep{kirkpatrick1983}.

Applied to redistricting bisection, the ``solution'' is a partition of $m$
tracts into two districts; a ``proposal'' is a boundary-tract flip
(reassign one tract from its current district to the other, maintaining
contiguity and balance); and the objective is the weighted edge cut $\EC$.
The SA loop is seeded from the \metis{} solution, so it starts from a
reasonable local optimum and refines further.

Three properties of this design are critical:

\begin{enumerate}
\item \textbf{METIS initialisation.}
  SA needs a good starting point.
  Starting from a random partition requires prohibitively many steps.
  Starting from \metis{}'s output means SA is refining a near-optimal
  solution rather than discovering one from scratch.

\item \textbf{Scaled step count.}
  A root node with $m = 2{,}500$ tracts is much harder to refine than a
  leaf node with $m = 50$ tracts.
  Setting $n_{\mathrm{steps}} = \texttt{steps\_per\_tract} \times m$ gives
  proportional computational effort across all tree levels.

\item \textbf{Best-ever tracking.}
  SA can increase $\EC$ during the high-temperature phase (exploration).
  Returning the final plan discards the minimum found mid-run.
  Tracking and returning the best-ever plan ensures exploration overhead
  is not wasted.
\end{enumerate}

\subsection{Main Results}

\paragraph{Algorithm.}
We define \sa{} bisection (Section~\ref{sec:algorithm}) as a drop-in
replacement for the \metis{} call at each bisection tree node.
The algorithm has two tunable parameters:
\texttt{steps\_per\_tract} $\in \{1,5,10,50\}$ controls runtime;
\texttt{t0\_factor} controls the initial acceptance probability.
A node-specific seed is derived deterministically from the global base
seed and the node's path in the bisection tree.

\paragraph{Acceptance probability analysis.}
For the default \texttt{t0\_factor}$= 0.01$, we show (Proposition~\ref{prop:acceptance})
that the initial acceptance probability for a $+1$ edge-cut move is
$\exp(-1/T_0) \approx \exp(-100/\EC_0)$.
For $\EC_0 = 100$, this is approximately $37\%$---enough to escape \metis{}'s
local optimum.
For $\EC_0 = 1{,}000$, $T_0 = 10$ and the acceptance probability is
$\exp(-0.1) \approx 90\%$---appropriate for the high-EC root node.

\paragraph{Empirical comparison.}
On NC ($k=14$), WI ($k=8$), and TX ($k=38$), \sa{} with
$\texttt{steps\_per\_tract} = 10$ achieves 10--18\% Polsby-Popper
improvement over baseline \metis{} (Section~\ref{sec:empirical}).
This lies between multi-seed \metis{} and \be{} at $T=100$.

\paragraph{Parameter sensitivity.}
Section~\ref{sec:sensitivity} shows that $\texttt{steps\_per\_tract} = 10$
is the practical sweet spot: steps$= 1$ is too short to escape the local
optimum; steps$= 50$ gives marginal further gain at $5\times$ runtime cost.
The \texttt{t0\_factor} is less critical; values in $[0.005, 0.05]$ all
perform similarly.

\subsection{Paper Structure}

Section~\ref{sec:background} reviews SA theory and the redistricting
structure-layer context.
Section~\ref{sec:algorithm} gives the formal SA bisection algorithm with
pseudocode, the node-seed derivation, and the acceptance probability
proposition.
Section~\ref{sec:empirical} presents empirical comparisons on NC, WI, and TX.
Section~\ref{sec:sensitivity} analyzes parameter sensitivity.
Section~\ref{sec:conclusion} situates \sa{} in the broader B-series
algorithm landscape.

\subsection{Relation to Prior Work}

\sa{} is a Structure-layer algorithm, orthogonal to the SeedCompositor
strategies (\cs{}, \be{}).
It is related to but distinct from \recom{}~\citep{deford2021}, which is
a Markov chain \emph{sampler} rather than an optimiser.
\citet{battiti1999} applied SA to graph partitioning but did not consider
the redistricting context or the METIS-initialisation design.
\geosec{}~\citep{b8paper} and \areasec{}~\citep{b9paper} are the closest
B-series relatives: all three modify what happens at each bisection node.
\sa{} differs in that it changes the \emph{solver} rather than the
\emph{objective}.
